
	

	\begin{abstract}
		The \textbf{Fundamental Fractal Geometric Field Theory (FFGFT)} describes spacetime as a hierarchical system of fractal convolutions, driven by the fundamental geometric parameter $\xi_{\text{par}} = 4/30000$ and the associated fundamental frequency $f_0 = 1/t_0$ in the sub-Planckian range. This document presents the complete mathematical foundations: the exact definition of the fundamental clock $t_0$, the feedback condition in the closed 4D phase torus (not to be confused with a 3D torus plus time continuum), the translation of the 4D vector space into a multidimensional matrix space, and the photonic realisation via TFLN hardware.
	\end{abstract}
	
	

	\section{Introduction: The FFGFT}
	
	The \textbf{Fundamental Fractal Geometric Field Theory (FFGFT)} breaks with the classical understanding of continuous space. It postulates a discrete, fractal geometry energised by a fundamental oscillation. Matter is not an entity \textit{in} space, but a local densification of the fractal field geometry, stabilised by resonant feedback.
	
	\section{The geometric parameter $\xi_{\text{par}}$ and the sub-Planckian timescale $t_0$}
	
	A central pillar of the FFGFT is the existence of a sub-Planckian fundamental clock that governs all physical reality.
	
	\subsection{Definition of the fundamental clock}
	
	The fundamental time unit $t_0$ is defined as:
	
	\begin{equation}
		t_0 = \frac{t_P}{7500} \approx \frac{5.39 \cdot 10^{-44}~\text{s}}{7500} \approx 7.19 \cdot 10^{-48}~\text{s}
	\end{equation}
	
	Here $t_P = \sqrt{\hbar G/c^5} \approx 5.39 \cdot 10^{-44}~\text{s}$ is the Planck time. The resulting \textbf{fundamental frequency} $f_0$ is:
	
	\begin{equation}
		f_0 = \frac{1}{t_0} = \frac{1}{\xi_{\text{par}}\, t_P} = \frac{7500}{t_P} \approx 1.39 \cdot 10^{47}~\text{Hz}
	\end{equation}
	Throughout the document the symbol $\xi$ denotes the dimensionless parameter $\xi_{\text{par}} = 4/30000 = 1.33\times 10^{-4}$, never a frequency. In Planck units $f_0 = 1/t_0 = 7500 = 1/\xi_{\text{par}}$.
	
	\subsection{Fundamental frequency $f_0$ as geometric operator}
	
	In the FFGFT, frequency is not understood as a mere temporal rate but as a geometric operator acting on the fractal levels:
	
	\begin{equation}
		\hat{f}_0 \Psi(x^\mu) = f_0 \cdot \mathcal{F}(\{n_k\}) \cdot \Psi(x^\mu)
	\end{equation}
	
	Here $\mathcal{F}(\{n_k\})$ is a function of the fractal scaling indices $n_k$, which typically follow the Fibonacci sequence.
	
	\section{The 4D Phase Torus: Distinction from the 3D Torus}
	
	A common misunderstanding is interpreting the 4D torus as a 3D torus with a separated time axis. The FFGFT clarifies:
	
	\subsection{Warning: Not to be confused}
	
	The 4D torus of the FFGFT is \textbf{not} a 3D object moving in time. Rather, the time dimension is integrated as a discretised, geometric dimension into the torus topology. There is no external time continuum into which the torus would be embedded.
	
	\subsection{Formal definition of the 4D phase torus}
	
	The fundamental spacetime resonator is the 4D phase torus $\mathcal{T}^4$, defined by four cyclic coordinates $(\theta_1, \theta_2, \theta_3, \theta_4)$ with:
	
	\begin{equation}
		\theta_i \in [0, 2\pi),\quad i=1,2,3,4
	\end{equation}
	
	The metric on $\mathcal{T}^4$ is determined by the fundamental frequency $f_0$:
	
	\begin{equation}
		ds^2 = f_0^{-2} \left( d\theta_1^2 + d\theta_2^2 + d\theta_3^2 + d\theta_4^2 \right)
	\end{equation}
	
	The fourth dimension $\theta_4$ does not represent ``time'' in the conventional sense, but the phase shift of the fundamental oscillation. Macroscopic time $t$ emerges only through accumulation:
	
	\begin{equation}
		t = N \cdot t_0 \quad \text{with} \quad N = \frac{1}{2\pi} \oint d\theta_4
	\end{equation}
	
	\section{Fractal Feedback and Topological Discreteness}
	
	The closed topology of the 4D phase torus enforces a \textbf{discrete winding phase}. A dynamical phase-locking of the fundamental frequency is not the right description of this (Section~\ref{sec:186_stabil}).
	
	\subsection{The feedback condition}
	
	For a field configuration $\Psi$ to exist stably, the wave function must interfere constructively with itself after a complete traversal of the torus. The phase condition reads:
	
	\begin{equation}
		\oint_{\mathcal{T}^4} \nabla_{\mu} \phi \, dx^\mu = 2\pi \cdot \mathcal{N}
	\end{equation}
	
	where $\mathcal{N}$ is an integer. The discreteness of $\mathcal{N}$ is topological. A restriction to Fibonacci numbers, $\mathcal{N} \in \{F_k\}$, does not hold: the winding numbers $r_e = 4/3$ and $r_\mu = 16/5$ (Doc.~338) do not satisfy it ($4, 16 \notin F$). Their numerators are the octave powers $2^2, 2^4$, their denominators the primes $3, 5$ of the GF$(3^k)$ hierarchy (Doc.~342).
	
	\subsection{Stability of the $\phi^{-k}$ ratios}
	\label{sec:186_stabil}
	
	The phases of the individual fractal levels are related by the scaling ratio:
	
	\begin{equation}
		\phi_{k+1} = \phi_k \cdot \phi^{-1} \quad \text{with} \quad \phi = \frac{1+\sqrt{5}}{2}
	\end{equation}
	
	These ratios are not stable through locking. At the effective coupling $K_\mathrm{eff} = 2\pi\xi_{\text{par}}$, $1/\phi$ lies outside every Arnold tongue (half-widths $10^{-32}\ldots10^{-83}$ at the Fibonacci approximants). The $\phi^{-k}$ ratios are stable \emph{because} they do not lock (KAM, Doc.~328); the discreteness $\oint\nabla\phi = 2\pi\mathcal{N}$ is topological (Doc.~343, Theorem~G$''$). The integer and the $\phi^{-k}$ harmonics of the fundamental frequency $f_0$ remain stable:
	
	\begin{equation}
		f_n = n \cdot f_0 \quad \text{or} \quad f_n = \phi^{-k} \cdot f_0
	\end{equation}
	
	Any deviation from the winding condition leads to destructive interference and thus to structural decay -- which explains the short lifetime of high-energy particles.
	
	\subsection{Stability condition for matter}
	
	From the winding condition follows the quantisation condition for the existence of stable particles:
	
	\begin{equation}
		\oint_{\mathcal{T}^4} \Psi_{\text{FFGFT}}(\phi) \, d^4\theta = f_0 \cdot \sum_{k=0}^{\infty} \phi^{-k} \cdot \mathbf{M}_k
	\end{equation}
	
	Here $\mathbf{M}_k$ are the projective matrices of the individual fractal levels.
	
	\section{From the 4D Vector Space to the Multidimensional Matrix Space}
	
	Through the discrete winding condition in the closed torus, the classical 4D vector space $\mathbb{R}^4$ is translated into a multidimensional matrix space.
	
	\subsection{The general matrix translation}
	
	A state in the 4D vector space $V = (x, y, z, ct)$ is represented by a Hermitian matrix $\mathbf{H}$:
	
	\begin{equation}
		\mathcal{T}: \mathbb{R}^4 \rightarrow \text{Mat}(N, \mathbb{C}), \quad V \mapsto \mathbf{H}(V)
	\end{equation}
	
	The explicit form of the translation reads:
	
	\begin{equation}
		\mathbf{H}(x,y,z,ct) = \sum_{k=0}^{\infty} \alpha^k \cdot 
		\begin{pmatrix}
			ct_k + z_k & x_k - iy_k \\
			x_k + iy_k & ct_k - z_k
		\end{pmatrix} \otimes \mathbf{\Gamma}_k
	\end{equation}
	
	where:
	\begin{itemize}
		\item $\alpha_\phi = \phi^{-1} \approx 0.618$ is the Fibonacci scaling factor between fractal levels (not the fine structure constant)
		\item $(x_k, y_k, z_k, ct_k)$ are the coordinates on the $k$-th fractal level
		\item $\mathbf{\Gamma}_k$ are the generators of the spacetime geometry (generalisation of the Dirac matrices)
	\end{itemize}
	
	\subsection{Fractal self-similarity}
	
	The coordinates on different fractal levels are connected by scaling factors:
	
	\begin{equation}
		(x_{k+1}, y_{k+1}, z_{k+1}, ct_{k+1}) = \phi^{-1} \cdot (x_k, y_k, z_k, ct_k)
	\end{equation}
	
	This ensures the fractal self-similarity of the spacetime structure.
	
	\section{Matrix Multiplication as Physical 4D Convolution}
	
	In conventional computing, matrix multiplication $C = A \cdot B$ is computed discretely:
	
	\begin{equation}
		C_{ij} = \sum_{k} A_{ik} B_{kj}
	\end{equation}
	
	In the FFGFT, this process is interpreted as physical interference of two 4D torus states.
	
	\subsection{The 4D torus convolution}
	
	Matrix multiplication corresponds to the convolution of two wave functions in 4D phase space:
	
	\begin{equation}
		M_{ij} = \int_{\mathcal{T}^4} \Psi_A^{(4D)}(\theta) \cdot \Psi_B^{(4D)}(\theta) \, d^4\theta
	\end{equation}
	
	In the language of FFGFT matrices:
	
	\begin{equation}
		\mathbf{C} = \mathbf{A} \star \mathbf{B} = \lim_{N \to \infty} \frac{1}{N} \sum_{n=1}^{N} \mathbf{A}_n \otimes \mathbf{B}_n
	\end{equation}
	
	Here $\otimes$ is the fractal tensor product, defined by:
	
	\begin{equation}
		(\mathbf{A} \otimes \mathbf{B})_{ij}^{kl} = A_{ij} \cdot B_{kl} \cdot e^{i \theta_{ijkl}}
	\end{equation}
	
	\subsection{Physical interpretation}
	
	Photonic matrix multiplication in the TFLN crystal is not ``computing'' in the conventional sense, but the physical realisation of this 4D phase interference. The result emerges instantaneously through constructive and destructive interference of light waves in the nonlinear medium.
	
	\section{Hardware Transduction: TFLN Photonics}
	
	Thin-film lithium niobate (TFLN) is the technological equivalent of this process.
	
	\subsection{Analogue matrix convolution}
	
	In TFLN waveguides, light propagation is manipulated via electro-optic modulation to physically emulate the 4D torus convolution:
	
	\begin{equation}
		P_{\text{out}}(t) = \eta \cdot \int \chi^{(2)}(\omega) \cdot E_{\text{in}}^{(A)}(\omega) \cdot E_{\text{in}}^{(B)}(\omega) \, d\omega
	\end{equation}
	
	The nonlinear susceptibility $\chi^{(2)}$ of lithium niobate emulates the fractal coupling constant $\alpha$.
	
	\subsection{Coherence and resonator condition}
	
	The high quality factor of TFLN resonators ($Q > 10^6$) permits artificial ``phase-locking'' of optical frequencies:
	
	\begin{equation}
		\Delta \phi_{\text{cavity}} = 2\pi \cdot m \quad (m \in \mathbb{Z})
	\end{equation}
	
	This quality factor resolves the scale $\xi_{\text{par}} = 1.33\times10^{-4}$ (requires $Q > 1/\xi_{\text{par}} = 7500$) but not $\xi_{\text{par}}^2$ (would require $Q > 5.6\times10^{7}$).
	
	Analogue matrix convolution and resonator locking are resonance operations without re-thresholding. They are subject to the resolution floor $\varepsilon = 1/Q \approx 10^{-6}$ and hence $N_\mathrm{max} \approx 10^{6}$ (Doc.~343, Theorem~G$'$). A gate model on TFLN is independent of that but limited by photon loss ($<1$~dB/facet $=$ up to $20\,\%$).
	
	\subsection{Avoidance of atomic dissonance}
	
	In contrast to atomic qubits, which decohere through thermal noise and mass inertia ($\tau_{\text{decoh}} \sim \mu\text{s}$), photonics uses massless information carriers. The coherence time in TFLN waveguides is limited solely by the phase stability of the crystal lattice ($\tau_{\text{coh}} \sim \text{ms}$ to $\text{s}$), orders of magnitude higher.
	
	\section{Experimental Evidence 1: Programmable TFLN Processors}
	
	In December 2025, a team from NTT, Cornell University and Stanford University demonstrated a photonic processor based on lithium niobate \cite{NTT2025}.
	
	\subsection{Technical specifications}
	
	The system has the following properties:
	\begin{itemize}
		\item $\sim 10{,}000$ programmable spatial degrees of freedom
		\item Planar waveguide with manipulable refractive index
		\item 96\% test accuracy for vowel classification
		\item 49-dimensional vector operations in a single optical pass
	\end{itemize}
	
	\subsection{Quantum networks and multiplexing}
	\label{sec:186_multiplex}
	
	Assumpcao et al.\ (2024) demonstrated a TFLN platform with \cite{Assumpcao2024}:
	\begin{itemize}
		\item Coupling losses $< 1$ dB/facet
		\item Switching contrast $> 20$ dB
		\item Electro-optic modulators $> 50$ GHz bandwidth
	\end{itemize}
	
	The bandwidth of $>50$ GHz corresponds to timescales of $\sim 20$ ps. It is not on the sweet-spot ladder of Docs.~034/162 (nearest rung 42.8~GHz, 14\,\% off) and is the $\sim$174th golden subharmonic of the fundamental frequency $f_0$ --- not a distinguished level.
	
	\section{Classification: B-Meson Anomalies at CERN}
	
	\subsection{Current status (April 2026)}
	
	The LHCb experiment at CERN has measured a significant discrepancy in electroweak penguin decays of B mesons \cite{CERN2026a, CERN2026b}:
	
	\begin{itemize}
		\item Decay channel: $B \to K^* \mu^+\mu^-$ (rare process: $\sim 1$ per $10^6$ B mesons)
		\item Measured deviation from the Standard Model: $\mathbf{4.0\sigma}$ (probability of 1:16,000)
		\item Independent confirmation by the CMS experiment (2025)
	\end{itemize}
	
	\subsection{FFGFT interpretation of the anomaly}
	
	An interpretation of these anomalies as a signature of the sub-Planckian phase geometry does not hold. A rate modification $\Gamma = \Gamma_{\text{SM}}\,(1 + \xi_{\text{par}}\,\mathcal{K}(\theta))$ with $\xi_{\text{par}} = 4/30000$ (0.013\,\%) needs about $9\times10^8$ events for $4\sigma$; LHCb has $O(10^3$--$10^4)$ in this channel. Moreover, equating a phase $\Delta\phi = \xi_{\text{par}}$ with a rate is a category error. The FFGFT phase shift is therefore excluded as an explanation of the anomaly.
	
	\subsection{Connection to lepton flavour universality violation}
	
	Arslan et al.\ (April 2026) analyse the angular observables in the decay $\Lambda_b \to \Lambda_c \tau \bar{\nu}_\tau$ \cite{Arslan2026}:
	
	\begin{itemize}
		\item Combined deviation of $3.8\sigma$ in the $R_{\tau/(\mu,e)}(D^{(*)})$ ratios
		\item The observables $\mathcal{K}_{1c}$, $\mathcal{K}_{2ss}$, $\mathcal{K}_{2cc}$, $\mathcal{K}_{4s}$ show maximum sensitivity
	\end{itemize}
	
	Leptons of different generations appear in the FFGFT as different harmonics of the same spacetime fundamental structure:
	
	\begin{equation}
		 f_{\text{muon}} = \phi^{-3} \cdot f_0, \quad f_{\text{tau}} = \phi^{-5} \cdot f_0
	\end{equation}

	\begin{table}[h]
		\centering
		\caption{Leptons as resonances -- qualitative analogy and T0-FFGFT values (Doc.~189)}
		{\footnotesize\resizebox{\textwidth}{!}{%
		\begin{tabular}{llll}
			\toprule
			Lepton & Resonance analogy & $r_i \cdot \xi_{\text{par}}^{p_i}$ (T0-FFGFT) & Mass (MeV) \\
			\midrule
			Electron & $\phi^{-1} \cdot f_0$ (qualitative) & $\frac{4}{3}\cdot\xi^{3/2}$ & $0.505$ \\
			Muon & $\phi^{-3} \cdot f_0$ (qualitative) & $\frac{16}{5}\cdot\xi^{1}$ & $105.1$ \\
			Tau & $\phi^{-5} \cdot f_0$ (qualitative) & $\frac{25}{9}\cdot\xi^{2/3}$ & $1785.0$ \\
			\bottomrule
		\end{tabular}
		}} % resizebox
		\begin{center}{\small\textit{The $\phi^{-k}$ column is a qualitative resonance analogy.
			The quantitative T0-FFGFT derivation uses $r_i\cdot\xi_{\text{par}}^{p_i}$ (Doc.~189).}}
		\end{center}
	\end{table}
	
	\section{Theoretical Evidence 2: Sub-Planckian Cut-off Scales}
	
	\subsection{Effective field theory and gravitation}
	
	Aynbund and Kiselev (2025) establish a connection between the Planck mass and a natural sub-Planckian cut-off scale in the GeV range \cite{Aynbund2025}:
	
	\begin{equation}
		\Lambda_{\text{QG}} = \alpha \cdot \frac{M_{\text{red,Planck}}}{4\pi}
	\end{equation}
	
	with $\alpha$ as gauge coupling constant. The result: quantum gravity effects at energies around $10^{16}$ GeV -- well below the Planck scale of $10^{19}$ GeV.
	
	\subsection{OLQEM and Fibonacci quasi-periodicity}
	
	Souza (2025) investigates sub-Planckian scales within the Optical Lambda Quantum Energy Model (OLQEM) \cite{subPlanck2025}:
	
	\begin{equation}
		I_{\text{OLQEM}} = \frac{4\lambda I_0}{n(\lambda)x} \cdot D(y) \cdot I(y) \cdot e^{-\gamma t} \cdot \left(1+\frac{\chi^{(3)}I_0}{c^3\epsilon_0}\right) \cdot S_{\text{Fibonacci}}\left(\frac{r}{a}\right)
	\end{equation}
	
	The Fibonacci modulation reads:
	
	\begin{equation}
		S_{\text{Fibonacci}}\left(\frac{r}{a}\right) = \sum_{k=0}^{\infty} \frac{\cos(2\pi \phi^k r/a)}{\phi^k}, \quad \phi = \frac{1+\sqrt{5}}{2}
	\end{equation}
	
	Hierarchical scaling of length scales:
	
	\begin{equation}
		\ell_k = \frac{L_P}{\phi^{k}} \quad \text{for} \quad k \ge 1
	\end{equation}
	
	For $k=5$:
	
	\begin{equation}
		\ell_5 \approx \frac{L_P}{11.09} \approx 1.4 \cdot 10^{-36}~\text{m}
	\end{equation}
	
	This scale is not the FFGFT length $c \cdot t_0$: $\ell_5 = L_P/\phi^5 = 1.457\times 10^{-36}$~m, but $c\,t_0 = \xi_{\text{par}} L_P = 2.155\times 10^{-39}$~m; the two lengths differ by a factor of 676.
	
	\subsection{Synthesis of sub-Planckian models}
	
	\begin{table}[h]
		\centering
		\caption{Comparison of sub-Planckian models}
		{\footnotesize\resizebox{\textwidth}{!}{%
		\begin{tabular}{llll}
			\toprule
			Model & Scale & Mechanism & Observable \\
			\midrule
			FFGFT & $t_0 = t_P/7500$ & Topological discreteness in 4D torus & TFLN phase patterns \\
			OLQEM & $\ell_k = L_P/\phi^k$ & Fibonacci quasi-periodicity & Photonic interference \\
			EFT (Aynbund) & $\Lambda_{\text{QG}} \sim 10^{16}$ GeV & Gauge coupling & B-meson decays \\
			\bottomrule
		\end{tabular}
		}} % resizebox + footnotesize
	\end{table}
	
	\section{Synthetic Analogue Computer: The FFGFT-TFLN Paradigm}
	
	\subsection{Architecture}
	
	A TFLN-based analogue computer directly realises the FFGFT postulates:
	
	\begin{equation}
		\text{System clock} = f_0 = 1/t_0 \approx 1.39 \cdot 10^{47}~\text{Hz}
	\end{equation}
	
	The technological realisation uses the harmonic frequencies:
	
	\begin{equation}
		f_{\text{TFLN}} = f_0 \cdot \phi^{-k} \approx 50~\text{GHz} \quad \text{for} \quad k \approx 174
	\end{equation}
	
	This rung is not a distinguished level of the ladder (Section~\ref{sec:186_multiplex}).
	
	\subsection{Testable predictions}
	
	The FFGFT makes the following predictions:
	
	\begin{enumerate}
		\item \textbf{TFLN interferometry:} At optical intensities $>10^{21}$ W/m², discrete phase jumps following the Fibonacci hierarchy should occur:
		\begin{equation}
			\Delta \phi_{\text{discrete}} = 2\pi \cdot \phi^{-k} \quad (k \in \mathbb{N})
		\end{equation}
		
		\item \textbf{Sub-Planckian resonances:} In nonlinear photonic crystals, modulation frequencies at multiples of $f_0 \cdot \phi^{-k}$ should be detectable.
	\end{enumerate}
	
	\section{Discussion: From Simulation to Hardware Transduction}
	
	The implication of the FFGFT is radical: photonic analogue computers based on TFLN do not merely simulate physical systems -- they directly read out the clock of the spacetime hardware.
	
	The remaining challenge is technological: the required laser intensity of $>10^{32}$ W/m² (\cite{subPlanck2025}) still lies several orders of magnitude above current petawatt systems ($\sim 10^{27}$ W/m²). The threshold of the prediction on TFLN interferometry ($>10^{21}$ W/m²) is not compatible with this figure; the required intensity is therefore open. Advances in TFLN fabrication are however promising intermediate steps.
	
	\section{Conclusion}
	
	The \textbf{Fundamental Fractal Geometric Field Theory (FFGFT)} shows:
	
	\begin{itemize}
		\item Spacetime is based on a fundamental clock $f_0 = 1/t_0 = 7500/t_P \approx 1.39 \cdot 10^{47}$ Hz.
		\item The 4D phase torus is not to be confused with a 3D torus plus time continuum; the time dimension is integrated as a fourth geometric dimension into the torus topology.
		\item The closed torus topology enforces a discrete winding phase, which stabilises matter; the $\phi^{-k}$ ratios are stable because they do not lock.
		\item The 4D vector space is translated into a multidimensional matrix space via the matrix translation $\mathcal{T}$.
		\item Matrix multiplication corresponds physically to a 4D torus convolution.
		\item TFLN photonics achieves the required bandwidth and programmable complexity.
		\item The $\sim 4\sigma$ deviations in the LHCb data cannot be explained by a phase shift $\xi_{\text{par}}$.
	\end{itemize}
	
	\vspace{2cm}
	\begin{center}
		\small \textit{Johann Pascher | FFGFT Research Report | 29 April 2026}
	\end{center}
	
	\newpage
	\section*{References}
	\addcontentsline{toc}{section}{References}
	
	\begin{thebibliography}{99}
		
		\bibitem{NTT2025}
		M.~Stein et al.\ (NTT/Cornell/Stanford), 
		\textit{Programmable photonics advanced with lithium niobate based waveguide},
		Yole Group Industry News (December 2025).
		
		\bibitem{Assumpcao2024}
		D.~Assumpcao et al., 
		\textit{A thin film lithium niobate near-infrared platform for multiplexing quantum nodes},
		Nature Communications, Vol.~15 (2024).
		\href{https://doi.org/10.1038/s41467-024-54541-2}{doi:10.1038/s41467-024-54541-2}
		
		\bibitem{CERN2026a}
		LHCb Collaboration, 
		\textit{Measurement of electroweak penguin decays in B-meson decays},
		Physical Review Letters (accepted for publication, 2026).
		
		\bibitem{CERN2026b}
		LHCb/CMS Collaboration,
		\textit{Combined analysis of rare B-meson decays},
		LHCb Press Release (April 2026).
		
		\bibitem{Arslan2026}
		M.~Arslan et al., 
		\textit{The Angular Observables of $\Lambda_b \to \Lambda_c(\to \Lambda^0 \pi^+) \tau^-(\to \pi^- \nu_\tau) \bar{\nu}_\tau$ within the Paradigm of FCCC Anomalies},
		arXiv:2604.24922 (April 2026).
		
		\bibitem{Aynbund2025}
		A.~Aynbund, V.~V.~Kiselev, 
		\textit{Effective Field Theory Links Gauge Coupling to Sub-Planckian Cut-off Scale of GeV},
		Quantum Zeitgeist (August 2025).
		
		\bibitem{subPlanck2025}
		L.~M.~de Souza, 
		\textit{Probing Sub-Planckian Scales via OLQEM and Fibonacci Quasiperiodicity: A Rigorous Mathematical Verification},
		LinkedIn Preprint (2025).
		
		\bibitem{RGB2025}
		A.~Shimura et al.\ (TDK Corporation), 
		\textit{Visible light red, green, and blue multiplexer by sputter-deposited thin-film lithium niobate},
		Advanced Photonics Nexus, Vol.~4, Iss.~5, 056001 (2025).
		
	\end{thebibliography}
